% amsart is the American Mathematical Society's article class. It is single-column and
% sets the title, the author addresses and the abstract in the AMS layout; many
% mathematics journals (and arXiv authors) use it. 11pt, a4paper or letterpaper are
% the usual options; "reqno" puts equation numbers on the right (the default here).
\documentclass[11pt,a4paper]{amsart}
\usepackage{amssymb}                 % amsart already loads amsmath and amsthm
\usepackage{mathtools}               % \coloneqq (the := sign, with correct spacing)
\usepackage[hidelinks]{hyperref}     % load last (or nearly); links in refs and citations

% Theorem-like environments. All share one counter, numbered within the section, so
% "Lemma 2.3" always comes after "Theorem 2.2": readers can find results by number.
\theoremstyle{plain}   % italic text: results
\newtheorem{theorem}{Theorem}[section]
\newtheorem{lemma}[theorem]{Lemma}
\newtheorem{proposition}[theorem]{Proposition}
\newtheorem{corollary}[theorem]{Corollary}
\theoremstyle{definition}  % upright text: definitions and conventions
\newtheorem{definition}[theorem]{Definition}
\newtheorem{example}[theorem]{Example}
\theoremstyle{remark}      % upright, run-in title, not numbered by default
\newtheorem{remark}[theorem]{Remark}
\numberwithin{equation}{section}  % equations as (2.1), (2.2): restart in each section

% Operators get their own commands so that the spacing is right.
\DeclareMathOperator{\diam}{diam}
\DeclareMathOperator{\Lip}{Lip}
\newcommand{\R}{\mathbb{R}}
\newcommand{\norm}[1]{\lVert #1\rVert}

\begin{document}

\title[Short running title]{A Fixed Point Approach to Existence\\ for Ordinary Differential Equations}

% Authors: \address, \email and \urladdr come after \author. Use one \author per person.
\author{First Author}
\address{Department of Mathematics, Example University, City, Country}
\email{first.author@example.edu}
\urladdr{https://example.edu/~first}

\author{Second Author}
\address{Institute of Analysis, Another University, City, Country}
\email{second.author@example.org}

% Funding and "received" notes go in \thanks; \dedicatory is optional.
\thanks{The first author was supported by Funding Body, grant 000000.}
\subjclass[2020]{Primary 34A12; Secondary 47H10, 54H25}  % https://mathscinet.ams.org/mathscinet/msc/msc2020.html
\keywords{fixed point, contraction mapping, initial value problem, Picard iteration}
\date{\today}

\begin{abstract}
State the question, the result and the idea of the proof in a few sentences. In
mathematics the abstract is read without the paper at hand: avoid citations and
displayed formulas, and say what the theorem is. Here we recover the local existence
and uniqueness of solutions of $x' = f(t,x)$ from the contraction mapping principle.
\end{abstract}

\maketitle

\section{Introduction}\label{sec:intro}
Say what is known, what is missing and what you prove, with references to
Theorem~\ref{thm:main} and to the literature~\cite{coddington55,rudin76}. Introduce the
notation in the order the reader needs it. Cite a specific place if you can,
for example \cite[Chapter~1]{coddington55}.

\section{Preliminaries}\label{sec:prelim}
\begin{definition}\label{def:contraction}
  Let $(X,d)$ be a metric space. A map $T\colon X\to X$ is a \emph{contraction} if there is
  $q\in[0,1)$ such that $d(Tx,Ty)\le q\,d(x,y)$ for all $x,y\in X$.
\end{definition}

\begin{lemma}[Banach \cite{banach22}]\label{lem:banach}
  If $(X,d)$ is complete and non-empty and $T$ is a contraction (Definition~\ref{def:contraction}) with constant $q$, then
  $T$ has exactly one fixed point $x^*$, and $d(T^nx_0,x^*)\le q^n(1-q)^{-1}d(Tx_0,x_0)$ for every $x_0\in X$.
\end{lemma}
\begin{proof}
  Put $x_n = T^n x_0$. By induction $d(x_{n+1},x_n)\le q^n d(x_1,x_0)$, so for $m>n$
  \[
    d(x_m,x_n)\le \sum_{k=n}^{m-1} q^k\, d(x_1,x_0)\le \frac{q^n}{1-q}\,d(x_1,x_0).
  \]
  Hence $(x_n)$ is Cauchy and converges to some $x^*$, which is fixed because $T$ is
  continuous. If $y$ is another fixed point, then $d(x^*,y)\le q\,d(x^*,y)$, so $y=x^*$.
  Letting $m\to\infty$ above gives the estimate.
\end{proof}

\begin{remark}\label{rem:completeness}
  Remarks and examples share the counter of the results, so ``Remark~\ref{rem:completeness}''
  is found in the same sequence. Completeness cannot be dropped in Lemma~\ref{lem:banach}.
\end{remark}

\section{An equation-heavy section}\label{sec:main}
Let $f\colon [0,a]\times\R^n\to\R^n$ be continuous and Lipschitz in the second variable,
\begin{equation}\label{eq:lipschitz}
  \norm{f(t,x)-f(t,y)}\le L\,\norm{x-y}.
\end{equation}
We write the initial value problem $x'=f(t,x)$, $x(0)=x_0$ as an integral equation, so
that a solution is a fixed point of
\begin{equation}\label{eq:picard}
  (Tx)(t) \coloneqq x_0 + \int_0^t f\bigl(s,x(s)\bigr)\,ds .
\end{equation}
Multi-line derivations use \texttt{align}; number only the lines you refer to, and let the
equals signs line up:
\begin{align}
  \norm{(Tx)(t)-(Ty)(t)}
    &\le \int_0^t \norm{f(s,x(s))-f(s,y(s))}\,ds \notag\\
    &\le L \int_0^t \norm{x(s)-y(s)}\,ds \label{eq:step}\\
    &\le L\,t\,\sup_{0\le s\le t}\norm{x(s)-y(s)}. \notag
\end{align}
Cases are typeset with the \texttt{cases} environment:
\begin{equation}\label{eq:weight}
  w_\lambda(t) =
  \begin{cases}
    e^{-\lambda t}, & 0\le t\le a,\\
    0, & \text{otherwise.}
  \end{cases}
\end{equation}

\begin{theorem}\label{thm:main}
  Under \eqref{eq:lipschitz} the problem $x'=f(t,x)$, $x(0)=x_0$ has exactly one solution on $[0,a]$.
\end{theorem}
\begin{proof}
  On $X=C([0,a];\R^n)$ use the norm $\norm{x}_\lambda=\sup_t e^{-\lambda t}\norm{x(t)}$ with
  $\lambda=2L$. It is equivalent to the supremum norm, so $X$ is complete. By \eqref{eq:step},
  \[
    e^{-\lambda t}\norm{(Tx)(t)-(Ty)(t)} \le L\,\norm{x-y}_\lambda\int_0^t e^{-\lambda(t-s)}\,ds
    \le \tfrac12\norm{x-y}_\lambda ,
  \]
  so $T$ is a contraction on $X$, and Lemma~\ref{lem:banach} gives the unique fixed point.
\end{proof}

\begin{corollary}
  The solution depends continuously on $x_0$: $\norm{x-y}_\lambda\le 2\norm{x_0-y_0}$.
\end{corollary}

\section{Conventions worth knowing}\label{sec:conv}
Put a full stop or comma after displayed formulas that end a sentence, as in
\eqref{eq:picard}. Use \verb|\eqref| for equations and plain \verb|\ref| for theorems, and
a non-breaking space (\verb|~|) before them. Write $\diam A$ and $\Lip(f)$ with the
operators declared in the preamble rather than as italic letters.

% The AMS ships two BibTeX styles: amsplain (numbered, as here) and amsalpha (labels such
% as [Rud76]). Change the style name to switch; the .bib file stays the same.
\bibliographystyle{amsplain}
\bibliography{refs}

\end{document}
